\section{AC-DC Converters (Rectifiers) \& Power Factor Correction (Week 8)}

\subsection{Keywords \& Comprehensive Overview}
\noindent\textbf{Ctrl+F keywords:} ac-dc rectifiers, uncontrolled rectifier, controlled rectifier, thyristor, SCR, firing angle, conduction angle, extinction angle, half-wave, full-wave, 1-phase bridge, 3-phase bridge, 3-pulse, 6-pulse, PIV, freewheeling diode, power triangle, real power, reactive power, apparent power, displacement power factor, DPF, distortion factor, DF, total harmonic distortion, THD, true power factor, PF, form factor, ripple factor, boost PFC, power factor correction, PFC inductor sizing, PFC bulk capacitor sizing, double-line frequency ripple, inverter mode, line-commutated inverter.

\vspace{0.5em}
\noindent AC-DC converters (rectifiers) convert alternating voltage into unidirectional direct voltage. Uncontrolled rectifiers employ power diodes whose commutation is dictated solely by circuit voltages, while controlled rectifiers utilize thyristors (SCRs) triggered at an adjustable delay firing angle $\alpha$. Passive diode rectifiers drawing pulsed currents produce severe line harmonics, motivating active Power Factor Correction (PFC) circuits to maintain near-unity power factor and satisfy regulatory limits (e.g., IEC 61000-3-2).

\subsubsection{Master Summary Table: AC-DC Converter Topologies}
\begin{center}
\footnotesize
\renewcommand{\arraystretch}{1.20}
\begin{tabularx}{\textwidth}{@{} >{\raggedright\arraybackslash}p{2.6cm} >{\raggedright\arraybackslash}p{2.5cm} >{\raggedright\arraybackslash}p{2.1cm} >{\raggedright\arraybackslash}X >{\centering\arraybackslash}p{1.7cm} >{\centering\arraybackslash}p{0.9cm} >{\centering\arraybackslash}p{0.9cm} >{\centering\arraybackslash}p{1.1cm} @{}}
\toprule
\textbf{Topology} & \textbf{DC Output $V_{dc}$} & \textbf{RMS Output $V_{rms}$} & \textbf{Controlled $V_{dc}(\alpha)$} & \textbf{PIV Rating} & \textbf{FF} & \textbf{RF} & $f_{\mathrm{rip}}$ \\
\midrule
\textbf{1-$\Phi$ Half-Wave} \newline (Diode, $R$ load) &
$\dfrac{V_m}{\pi} \approx 0.318 V_m$ \newline $= 0.450 V_s$ &
$\dfrac{V_m}{2} = 0.500 V_m$ \newline $= 0.707 V_s$ &
$\dfrac{V_m}{2\pi}(1+\cos\alpha)$ \newline ($R$ load) &
$V_m = \sqrt{2} V_s$ &
$1.571$ &
$1.211$ &
$f_{\mathrm{line}}$ \\
\midrule
\textbf{1-$\Phi$ Full-Bridge} \newline (Diode / SCR) &
$\dfrac{2V_m}{\pi} \approx 0.637 V_m$ \newline $= 0.900 V_s$ &
$\dfrac{V_m}{\sqrt{2}} \approx 0.707 V_m$ \newline $= V_s$ &
$\dfrac{2V_m}{\pi}\cos\alpha$ (cont. $L$) \newline
$\dfrac{V_m}{\pi}(1+\cos\alpha)$ ($R$) &
$V_m = \sqrt{2} V_s$ &
$1.111$ &
$0.483$ &
$2 f_{\mathrm{line}}$ \\
\midrule
\textbf{1-$\Phi$ Center-Tap} \newline (Full-wave) &
$\dfrac{2V_m}{\pi} \approx 0.637 V_m$ &
$\dfrac{V_m}{\sqrt{2}} \approx 0.707 V_m$ &
$\dfrac{2V_m}{\pi}\cos\alpha$ (cont. $L$) &
$2V_m = 2\sqrt{2} V_s$ &
$1.111$ &
$0.483$ &
$2 f_{\mathrm{line}}$ \\
\midrule
\textbf{3-$\Phi$ Half-Wave} \newline (3-pulse, per-phase) &
$\dfrac{3\sqrt{3}V_{m,ph}}{2\pi} \approx 0.827 V_{m,ph}$ \newline $= 1.170 V_{ph,rms}$ &
$0.841 V_{m,ph}$ \newline $= 1.190 V_{ph,rms}$ &
$\dfrac{3\sqrt{3}V_{m,ph}}{2\pi}\cos\alpha$ \newline (continuous current) &
$\sqrt{3}V_{m,ph} = V_{m,LL}$ \newline $= \sqrt{2} V_{LL,rms}$ &
$1.017$ &
$0.183$ &
$3 f_{\mathrm{line}}$ \\
\midrule
\textbf{3-$\Phi$ Full-Bridge} \newline (6-pulse, line-line) &
$\dfrac{3V_{m,LL}}{\pi} \approx 1.350 V_{LL,rms}$ \newline $= 2.339 V_{ph,rms}$ &
$1.352 V_{LL,rms}$ \newline $= 2.341 V_{ph,rms}$ &
$\dfrac{3V_{m,LL}}{\pi}\cos\alpha$ \newline $= 1.350 V_{LL,rms}\cos\alpha$ &
$V_{m,LL} = \sqrt{2} V_{LL,rms}$ \newline $= \sqrt{6} V_{ph,rms}$ &
$1.001$ &
$0.042$ &
$6 f_{\mathrm{line}}$ \\
\bottomrule
\end{tabularx}
\end{center}

---

\subsection{AC Power Fundamentals, Harmonics \& Power Factor}

\subsubsection{Instantaneous, Real, Reactive, and Apparent Power}
Consider a single-phase AC voltage source supplying a load with instantaneous voltage and current:
\begin{align}
v(t) &= \sqrt{2} V_s \sin(\omega t) = V_m \sin(\omega t) \\
i(t) &= \sqrt{2} I_s \sin(\omega t - \phi) = I_m \sin(\omega t - \phi)
\end{align}
where $V_s, I_s$ are the RMS values, and $\phi$ is the phase displacement angle of current relative to voltage ($\phi > 0$ for lagging/inductive load; $\phi < 0$ for leading/capacitive load).

\paragraph{1. Instantaneous Power $p(t)$:}
\begin{equation}
p(t) = v(t) i(t) = 2 V_s I_s \sin(\omega t)\sin(\omega t - \phi) = \underbrace{V_s I_s \cos\phi \left[1 - \cos(2\omega t)\right]}_{\text{Unidirectional (active power transfer)}} - \underbrace{V_s I_s \sin\phi \sin(2\omega t)}_{\text{Alternating (reactive energy exchange)}}
\end{equation}

\paragraph{2. Active (Real) Power $P$:} The time-average of $p(t)$ over one line period $T = 2\pi/\omega$:
\begin{equation}
P = \frac{1}{T} \int_0^T p(t)\,dt = V_s I_s \cos\phi \quad [\si{\watt}\text{ or W}]
\end{equation}

\paragraph{3. Reactive Power $Q$:} The peak amplitude of the zero-average oscillating power:
\begin{equation}
Q = V_s I_s \sin\phi \quad [\si{\text{var}}\text{ or Volt-Ampere Reactive}]
\end{equation}

\paragraph{4. Apparent Power $S$:} The product of total RMS voltage and total RMS current:
\begin{equation}
S = V_s I_s = \sqrt{P^2 + Q^2} \quad [\si{\volt\ampere}\text{ or VA}]
\end{equation}

\subsubsection{The Power Triangle}
The geometric relationship between $P, Q,$ and $S$ forms the \textbf{Power Triangle}:
\begin{itemize}
  \item Horizontal adjacent leg: Real power $P = S \cos\phi$ (W).
  \item Vertical opposite leg: Reactive power $Q = S \sin\phi$ (var) ($+Q$ for lagging/inductive; $-Q$ for leading/capacitive).
  \item Hypotenuse: Apparent power $S = \sqrt{P^2 + Q^2}$ (VA).
  \item Displacement angle: $\phi = \arctan\left(\dfrac{Q}{P}\right)$.
\end{itemize}
\textit{Power Factor Correction via Parallel Capacitors:} When an inductive motor load draws lagging reactive power $Q_L$, placing a capacitor bank in parallel drawing leading reactive power $Q_C = -V_s^2 \omega C$ cancels the net reactive power ($Q_{\mathrm{net}} = Q_L - |Q_C|$), reducing total line current $I_s = \sqrt{P^2 + Q_{\mathrm{net}}^2}/V_s$ without changing the useful work $P$.

\subsubsection{Non-Sinusoidal Currents and Harmonic Fourier Decomposition}
Electronic and power electronic loads (rectifiers, SMPS, motor drives) draw non-sinusoidal periodic currents. Any periodic function of angular frequency $\omega = 2\pi/T$ satisfies the Fourier series expansion:
\begin{equation}
i_s(t) = I_{dc} + \sum_{n=1}^\infty \left[ a_n \cos(n\omega t) + b_n \sin(n\omega t) \right] = I_{dc} + \sum_{n=1}^\infty \sqrt{2} I_{sn} \sin(n\omega t - \phi_n)
\end{equation}
where:
\begin{equation}
a_n = \frac{2}{T}\int_0^T i_s(t)\cos(n\omega t)\,dt, \qquad b_n = \frac{2}{T}\int_0^T i_s(t)\sin(n\omega t)\,dt
\end{equation}
\paragraph{Waveform Symmetries:}
\begin{itemize}
  \item \textbf{Even symmetry} ($i(-t) = i(t)$): $b_n = 0$; only cosine terms exist.
  \item \textbf{Odd symmetry} ($i(-t) = -i(t)$): $I_{dc} = 0, a_n = 0$; only sine terms exist.
  \item \textbf{Half-wave symmetry} ($i(t + T/2) = -i(t)$): $I_{dc} = 0$, and all even harmonics vanish ($a_n = b_n = 0$ for $n = 2, 4, 6, \dots$). The waveform contains \emph{only odd harmonics} ($n = 1, 3, 5, 7, \dots$).
\end{itemize}

\paragraph{Active Power under Pure Sinusoidal Supply Voltage:}
Assuming a pure sinusoidal grid supply $v_s(t) = \sqrt{2} V_s \sin(\omega t)$:
\begin{equation}
P = \frac{1}{T}\int_0^T v_s(t) i_s(t)\,dt = V_s I_{s1} \cos\phi_1
\end{equation}
\textbf{Crucial Insight:} Harmonic currents ($n \ge 2$) interacting with the fundamental sinusoidal voltage produce \emph{zero net average active power} over a complete cycle ($\int_0^T \sin(\omega t)\sin(n\omega t - \phi_n)\,dt = 0$ for $n \ne 1$). However, they increase the total RMS current, causing extra $I^2 R$ transmission losses and heating in transformers, cables, and circuit breakers!

\subsubsection{Key Power Quality Metrics: DPF, DF, THD, and True PF}
When current contains harmonics, the total RMS current is:
\begin{equation}
I_{s,\mathrm{rms}} = \sqrt{I_{dc}^2 + \sum_{n=1}^\infty I_{sn,\mathrm{rms}}^2} = \sqrt{I_{s1,\mathrm{rms}}^2 + \sum_{n=2}^\infty I_{sn,\mathrm{rms}}^2}
\end{equation}

\paragraph{1. Displacement Power Factor (DPF):}
The cosine of the phase displacement angle $\phi_1$ between the fundamental voltage and fundamental current:
\begin{equation}
DPF = \cos\phi_1
\end{equation}

\paragraph{2. Distortion Factor (DF):}
The ratio of the fundamental RMS current to the total RMS current:
\begin{equation}
DF = \frac{I_{s1,\mathrm{rms}}}{I_{s,\mathrm{rms}}} = \frac{I_{s1}}{\sqrt{I_{s1}^2 + \sum_{n=2}^\infty I_{sn}^2}} \quad (0 \le DF \le 1)
\end{equation}

\paragraph{3. Total Harmonic Distortion (THD):}
The ratio of the RMS value of all harmonic components combined to the RMS value of the fundamental:
\begin{equation}
THD = \frac{\sqrt{\sum_{n=2}^\infty I_{sn,\mathrm{rms}}^2}}{I_{s1,\mathrm{rms}}} = \frac{\sqrt{I_{s,\mathrm{rms}}^2 - I_{s1,\mathrm{rms}}^2}}{I_{s1,\mathrm{rms}}} = \sqrt{\left(\frac{I_{s,\mathrm{rms}}}{I_{s1,\mathrm{rms}}}\right)^2 - 1} = \sqrt{\left(\frac{1}{DF}\right)^2 - 1}
\end{equation}
Inverting this expression directly yields the exact relation between $DF$ and $THD$:
\begin{equation}
DF = \frac{1}{\sqrt{1 + THD^2}}
\end{equation}

\paragraph{4. True (Total) Power Factor (PF):}
The fundamental definition of power factor is the ratio of active power $P$ to apparent power $S$:
\begin{equation}
PF = \frac{P}{S} = \frac{V_s I_{s1,\mathrm{rms}} \cos\phi_1}{V_s I_{s,\mathrm{rms}}} = \left(\frac{I_{s1,\mathrm{rms}}}{I_{s,\mathrm{rms}}}\right) \cos\phi_1 = DF \times DPF = \frac{\cos\phi_1}{\sqrt{1 + THD^2}}
\end{equation}
\begin{center}
\fbox{\parbox{0.92\textwidth}{
\textbf{The Dual Origin of Low Power Factor:}
\begin{enumerate}
  \item \textbf{Displacement:} Fundamental current lags/leads voltage ($\phi_1 \ne 0 \implies DPF < 1$).
  \item \textbf{Distortion:} Non-sinusoidal current pulses create harmonics ($THD > 0 \implies DF < 1$).
\end{enumerate}
Even if the current is perfectly in phase with voltage ($\phi_1 = 0^\circ \implies DPF = 1.0$), severe distortion ($THD = 100\%$) yields a true power factor of only $PF = 1/\sqrt{1 + 1^2} = 0.707$!
}}
\end{center}

\subsubsection{Waveform Quality Metrics: Form Factor ($FF$) and Ripple Factor ($RF$)}
For rectified DC voltages and currents:
\begin{align}
\textbf{Form Factor (FF):} \quad & FF = \frac{V_{rms}}{V_{dc}} \ge 1 \quad \text{($FF = 1.0$ for pure, ripple-free DC)} \\
\textbf{Ripple Factor (RF):} \quad & RF = \frac{V_{ac,rms}}{V_{dc}} = \frac{\sqrt{V_{rms}^2 - V_{dc}^2}}{V_{dc}} = \sqrt{\left(\frac{V_{rms}}{V_{dc}}\right)^2 - 1} = \sqrt{FF^2 - 1}
\end{align}
\textit{Efficiency of Rectification:} $\eta = \dfrac{P_{dc}}{P_{in}} = \dfrac{V_{dc} I_{dc}}{V_{rms} I_{rms}} = \dfrac{1}{FF_v \cdot FF_i}$. For a pure resistive load, $\eta = \dfrac{1}{FF^2}$.

\subsubsection{Example: Harmonic Spectrum and THD of an Ideal Square Wave}
A single-phase full-bridge rectifier with large smoothing inductor draws a square-wave line current of amplitude $I_d$ (from Week 8 Slides 38--43):
\begin{equation}
i_s(t) = \begin{cases} +I_d, & 0 \le \omega t < \pi \\ -I_d, & \pi \le \omega t < 2\pi \end{cases}
\end{equation}
\paragraph{1. Total RMS Current:}
\begin{equation}
I_{s,\mathrm{rms}} = \sqrt{\frac{1}{2\pi}\int_0^{2\pi} (\pm I_d)^2 \, d(\omega t)} = I_d
\end{equation}
\paragraph{2. Fourier Series Expansion:}
Because $i_s(t)$ has odd and half-wave symmetry, $a_n = 0$, and for odd $n$:
\begin{equation}
b_n = \frac{2}{\pi}\int_0^\pi I_d \sin(n\omega t)\,d(\omega t) = \frac{2 I_d}{n\pi}\left[-\cos(n\omega t)\right]_0^\pi = \frac{4 I_d}{n\pi} \quad (n = 1, 3, 5, \dots)
\end{equation}
\begin{equation}
i_s(t) = \frac{4 I_d}{\pi}\left[ \sin(\omega t) + \frac{1}{3}\sin(3\omega t) + \frac{1}{5}\sin(5\omega t) + \frac{1}{7}\sin(7\omega t) + \cdots \right]
\end{equation}
\paragraph{3. Fundamental RMS and Distortion Factor:}
\begin{equation}
I_{s1,\mathrm{rms}} = \frac{b_1}{\sqrt{2}} = \frac{4 I_d}{\sqrt{2}\pi} = \frac{2\sqrt{2}}{\pi} I_d \approx 0.9003 I_d
\end{equation}
\begin{equation}
DF = \frac{I_{s1,\mathrm{rms}}}{I_{s,\mathrm{rms}}} = \frac{2\sqrt{2}}{\pi} \approx 0.9003
\end{equation}
\paragraph{4. Total Harmonic Distortion (THD):}
\begin{equation}
THD = \sqrt{\left(\frac{1}{DF}\right)^2 - 1} = \sqrt{\left(\frac{\pi}{2\sqrt{2}}\right)^2 - 1} = \sqrt{\frac{\pi^2}{8} - 1} \approx \sqrt{1.2337 - 1} = \sqrt{0.2337} \approx 0.4834 \quad (48.34\%)
\end{equation}
If the bridge is uncontrolled ($\alpha = 0$), $DPF = \cos 0^\circ = 1.0$, giving a true power factor:
\begin{equation}
PF = DF \times DPF = 0.9003 \times 1.0 = 0.9003
\end{equation}

---

\subsection{Single-Phase Rectifier Topologies}

\subsubsection{Single-Phase Half-Wave Uncontrolled (Diode) Rectifier}
Supply voltage: $v_s(t) = V_m \sin(\omega t) = \sqrt{2} V_s \sin(\omega t)$.

\paragraph{1. Pure Resistive Load ($R$):}
Diode conducts during the positive half-cycle ($0 \le \omega t \le \pi$), blocked during negative half-cycle ($\pi \le \omega t \le 2\pi$).
\begin{align}
V_{dc} &= \frac{1}{2\pi}\int_0^\pi V_m \sin(\omega t)\,d(\omega t) = \frac{V_m}{2\pi}\left[-\cos(\omega t)\right]_0^\pi = \frac{V_m}{\pi} \approx 0.3183 V_m = \frac{\sqrt{2} V_s}{\pi} \approx 0.4502 V_s \\
V_{rms} &= \sqrt{\frac{1}{2\pi}\int_0^\pi V_m^2 \sin^2(\omega t)\,d(\omega t)} = \frac{V_m}{2} = 0.500 V_m = \frac{V_s}{\sqrt{2}} \approx 0.7071 V_s \\
FF &= \frac{V_{rms}}{V_{dc}} = \frac{V_m/2}{V_m/\pi} = \frac{\pi}{2} \approx 1.5708 \\
RF &= \sqrt{FF^2 - 1} = \sqrt{\left(\frac{\pi}{2}\right)^2 - 1} \approx 1.2114 \quad (121.1\%) \\
\text{PIV} &= V_m = \sqrt{2} V_s \quad (\text{Peak Inverse Voltage across diode})
\end{align}
\textit{Shortcomings:} High ripple factor ($121\%$), poor transformer utilization ($TUF \approx 0.286$), and net DC current in the AC source causing magnetic core saturation.

\paragraph{2. Inductive-Resistive Load ($R-L$ Load) and Extinction Angle $\beta$:}
When the load has inductance $L$, the circuit differential equation is:
\begin{equation}
L \frac{di}{dt} + R i = V_m \sin(\omega t)
\end{equation}
The stored inductive energy $\tfrac{1}{2} L i^2$ maintains current flow even after $v_s(t)$ reverses at $\omega t = \pi$. The diode remains forward-biased until current drops to zero at the \textbf{extinction angle} $\beta > \pi$:
\begin{equation}
i(\omega t) = \frac{V_m}{Z}\left[\sin(\omega t - \phi) + \sin\phi \cdot e^{-\frac{\omega t}{\tan\phi}}\right] \quad \text{for } 0 \le \omega t \le \beta
\end{equation}
where $Z = \sqrt{R^2 + (\omega L)^2}$, $\tan\phi = \frac{\omega L}{R}$. The extinction angle $\beta$ is the transcendental solution of $i(\beta) = 0$:
\begin{equation}
\sin(\beta - \phi) + \sin\phi \cdot e^{-\frac{\beta}{\tan\phi}} = 0
\end{equation}
During the interval $\pi \le \omega t \le \beta$, $v_d(t) = v_s(t) < 0$. This negative voltage \emph{reduces} the average DC output voltage:
\begin{equation}
V_{dc} = \frac{1}{2\pi}\int_0^\beta V_m \sin(\omega t)\,d(\omega t) = \frac{V_m}{2\pi}(1 - \cos\beta)
\end{equation}

\paragraph{3. Freewheeling Diode (FWD):}
Placing a diode $D_{fw}$ in anti-parallel across the $R-L$ load clamps the load voltage to zero when $v_s(t)$ goes negative at $\omega t = \pi$:
\begin{itemize}
  \item For $0 \le \omega t \le \pi$: Supply diode conducts; $v_d = v_s$.
  \item For $\pi \le \omega t \le 2\pi$: $D_{fw}$ turns ON, clamping $v_d = 0$. Stored magnetic energy circulates through $D_{fw}$ and $R$, decaying exponentially.
  \item Average output voltage returns to the resistive value: $V_{dc} = \dfrac{V_m}{\pi}$.
\end{itemize}

\subsubsection{Single-Phase Half-Wave Controlled (SCR) Rectifier}
Replacing the diode with a thyristor (SCR) allows phase-delay control via gate firing angle $\alpha$.

\paragraph{1. Resistive Load ($R$):}
Gate trigger pulse applied at $\omega t = \alpha$ ($0 \le \alpha \le \pi$). The SCR conducts from $\alpha$ until natural zero-crossing at $\pi$.
\begin{equation}
\text{Conduction angle:} \quad \theta = \pi - \alpha \quad (\alpha + \theta = 180^\circ)
\end{equation}
\begin{align}
V_{dc} &= \frac{1}{2\pi}\int_\alpha^\pi V_m \sin(\omega t)\,d(\omega t) = \frac{V_m}{2\pi}\left[-\cos(\omega t)\right]_\alpha^\pi = \frac{V_m}{2\pi}(1 + \cos\alpha) = \frac{\sqrt{2} V_s}{2\pi}(1 + \cos\alpha) \\
&= \frac{0.4502 V_s}{2}(1 + \cos\alpha) \approx 0.2251 V_s (1 + \cos\alpha) \\
I_{dc} &= \frac{V_{dc}}{R} = \frac{V_m}{2\pi R}(1 + \cos\alpha) \\
V_{rms} &= \sqrt{\frac{1}{2\pi}\int_\alpha^\pi V_m^2 \sin^2(\omega t)\,d(\omega t)} = \frac{V_m}{2}\sqrt{1 - \frac{\alpha}{\pi} + \frac{\sin(2\alpha)}{2\pi}} \\
I_{T,\mathrm{rms}} &= \frac{V_{rms}}{R} = \frac{V_m}{2R}\sqrt{1 - \frac{\alpha}{\pi} + \frac{\sin(2\alpha)}{2\pi}}
\end{align}
When $\alpha = 0^\circ$: $V_{dc} = \dfrac{V_m}{\pi}$. When $\alpha = 180^\circ$: $V_{dc} = 0$.

\paragraph{2. Firing Delay Time:}
The time delay $t_\alpha$ corresponding to firing angle $\alpha$ (in radians):
\begin{equation}
t_\alpha = \frac{\alpha}{\omega} = \frac{\alpha}{2\pi f_{\mathrm{line}}} \quad [\si{\second}]
\end{equation}

\subsubsection{Single-Phase Full-Wave Uncontrolled Diode Bridge Rectifier}
A four-diode bridge conducts on both half-cycles: $(D_1, D_2)$ during $v_s > 0$, and $(D_3, D_4)$ during $v_s < 0$.

\paragraph{1. Resistive Load ($R$):}
\begin{align}
V_{dc} &= \frac{1}{\pi}\int_0^\pi V_m \sin(\omega t)\,d(\omega t) = \frac{2 V_m}{\pi} \approx 0.6366 V_m = \frac{2\sqrt{2} V_s}{\pi} \approx 0.9003 V_s \\
V_{rms} &= \sqrt{\frac{1}{\pi}\int_0^\pi V_m^2 \sin^2(\omega t)\,d(\omega t)} = \frac{V_m}{\sqrt{2}} = V_s \approx 0.7071 V_m \\
FF &= \frac{V_{rms}}{V_{dc}} = \frac{V_s}{\frac{2\sqrt{2}}{\pi} V_s} = \frac{\pi}{2\sqrt{2}} \approx 1.1107 \\
RF &= \sqrt{FF^2 - 1} = \sqrt{\left(\frac{\pi}{2\sqrt{2}}\right)^2 - 1} \approx 0.4834 \quad (48.34\%) \\
\text{PIV} &= V_m = \sqrt{2} V_s \quad (\text{For center-tapped transformer: } \text{PIV} = 2V_m) \\
f_{\mathrm{ripple}} &= 2 f_{\mathrm{line}} \quad (100\text{ Hz for 50 Hz grid; } 120\text{ Hz for 60 Hz grid})
\end{align}

\paragraph{2. Capacitor-Filtered Diode Bridge (Peak Rectifier):}
In practical DC power supplies, a large electrolytic capacitor $C$ is connected across the load:
\begin{itemize}
  \item The output voltage is held near peak $V_m$ with ripple $\Delta V_o \approx \dfrac{I_{dc}}{2 f_{\mathrm{line}} C}$.
  \item The diodes conduct only for a narrow conduction interval $\Delta t_{\mathrm{on}} \ll T/2$ near the AC voltage peak when $v_s(t) > v_C(t)$.
  \item Consequently, the input current consists of narrow, high-amplitude pulses with $THD > 80\%$ and $PF \approx 0.55 - 0.65$.
\end{itemize}

\subsubsection{Single-Phase Full-Wave Controlled (SCR) Bridge Rectifier}
Four thyristors $(VT_1, VT_4)$ and $(VT_2, VT_3)$ form diagonal pairs triggered at $\alpha$ and $\pi + \alpha$.

\paragraph{1. Resistive Load ($R$):}
Current drops to zero at $\pi$, so thyristors extinguish at the end of each half-cycle:
\begin{align}
V_{dc} &= \frac{1}{\pi}\int_\alpha^\pi V_m \sin(\omega t)\,d(\omega t) = \frac{V_m}{\pi}(1 + \cos\alpha) = \frac{\sqrt{2} V_s}{\pi}(1 + \cos\alpha) \approx 0.4502 V_s (1 + \cos\alpha) \\
I_{dc} &= \frac{V_{dc}}{R} = \frac{\sqrt{2} V_s}{\pi R}(1 + \cos\alpha) \\
I_{T,\mathrm{avg}} &= \frac{I_{dc}}{2} = \frac{\sqrt{2} V_s}{2\pi R}(1 + \cos\alpha) \approx \frac{0.2251 V_s}{R}(1 + \cos\alpha) \\
V_{rms} &= \frac{V_m}{\sqrt{2}}\sqrt{1 - \frac{\alpha}{\pi} + \frac{\sin(2\alpha)}{2\pi}} \\
I_{T,\mathrm{rms}} &= \frac{V_m}{2R}\sqrt{1 - \frac{\alpha}{\pi} + \frac{\sin(2\alpha)}{2\pi}}
\end{align}

\paragraph{2. Highly Inductive Load (Continuous Conduction Mode, $L \to \infty$):}
When a large smoothing inductor $L_d$ ensures ripple-free DC current $i_d(t) = I_d$:
\begin{itemize}
  \item At $\omega t = \pi$, supply voltage $v_s$ reverses, but $VT_1, VT_4$ cannot turn off because the inductor forces current to continue flowing.
  \item The outgoing pair $VT_1, VT_4$ conducts with negative $v_d$ until the incoming pair $VT_2, VT_3$ is fired at $\omega t = \pi + \alpha$, which reverse-biases and commutates $VT_1, VT_4$.
  \item Conduction interval for each thyristor pair is exactly $\pi$ radians ($180^\circ$).
\end{itemize}
\begin{equation}
V_{dc} = \frac{1}{\pi}\int_\alpha^{\pi+\alpha} V_m \sin(\omega t)\,d(\omega t) = \frac{V_m}{\pi}\left[-\cos(\omega t)\right]_\alpha^{\pi+\alpha} = \frac{2 V_m}{\pi}\cos\alpha = \frac{2\sqrt{2} V_s}{\pi}\cos\alpha \approx 0.9003 V_s \cos\alpha
\end{equation}
\begin{align}
\text{Thyristor average current:} \quad & I_{T,\mathrm{avg}} = \frac{I_d}{2} = 0.500 I_d \\
\text{Thyristor RMS current:} \quad & I_{T,\mathrm{rms}} = \frac{I_d}{\sqrt{2}} \approx 0.7071 I_d \\
\text{Source RMS current:} \quad & I_{s,\mathrm{rms}} = I_d \\
\text{Source fundamental RMS:} \quad & I_{s1,\mathrm{rms}} = \frac{2\sqrt{2}}{\pi} I_d \approx 0.9003 I_d \\
\text{Distortion Factor:} \quad & DF = \frac{I_{s1,\mathrm{rms}}}{I_{s,\mathrm{rms}}} = \frac{2\sqrt{2}}{\pi} \approx 0.9003 \\
\text{Displacement angle:} \quad & \phi_1 = \alpha \implies DPF = \cos\alpha \\
\text{True Power Factor:} \quad & PF = DF \times DPF = \frac{2\sqrt{2}}{\pi}\cos\alpha \approx 0.9003 \cos\alpha
\end{align}

\subsubsection{Line-Commutated Inversion Mode ($\alpha > 90^\circ$)}
Examining $V_{dc} = \dfrac{2 V_m}{\pi}\cos\alpha$:
\begin{itemize}
  \item For $0^\circ \le \alpha < 90^\circ$: $\cos\alpha > 0 \implies V_{dc} > 0$. Power $P = V_{dc} I_d > 0$ flows from AC source to DC load (\textbf{Rectifier Mode}).
  \item At $\alpha = 90^\circ$: $\cos 90^\circ = 0 \implies V_{dc} = 0$.
  \item For $90^\circ < \alpha < 180^\circ$: $\cos\alpha < 0 \implies V_{dc} < 0$. Because thyristors are unidirectional switches ($I_d > 0$ strictly), the power is negative:
  \begin{equation}
  P_{dc} = V_{dc} I_d < 0 \quad \implies \quad \text{Power flows from DC side into AC grid (\textbf{Inverter Mode})}
  \end{equation}
\end{itemize}
\paragraph{Mandatory Physical Requirements for Inverter Operation:}
\begin{enumerate}
  \item An active DC voltage source / battery / motor back-EMF $E$ must be present on the DC side with polarity opposing the rectifier forward polarity ($E$ aiding current flow).
  \item Loop equation: $E + V_{dc} = I_d R \implies I_d = \dfrac{E - |V_{dc}|}{R} > 0 \implies E > |V_{dc}|$.
  \item Sufficient DC inductance $L_d$ to maintain continuous current.
  \item \textbf{Commutation Margin Angle $\gamma$:} The maximum firing angle is strictly limited:
  \begin{equation}
  \alpha_{\max} = 180^\circ - \gamma - \mu
  \end{equation}
  where $\mu$ is the overlap angle due to AC source inductance, and $\gamma$ is the thyristor turn-off margin angle ($\gamma = \omega t_q$, where $t_q$ is the thyristor turn-off recovery time, typically $15^\circ - 30^\circ$). Firing beyond $\alpha_{\max}$ causes \textbf{commutation failure}, resulting in an instantaneous DC short-circuit!
\end{enumerate}

---

\subsection{Three-Phase Rectifier Topologies}

\subsubsection{Three-Phase Half-Wave Rectifier (3-Pulse)}
A star-connected three-phase supply with phase voltages $v_{an}, v_{bn}, v_{cn}$ of peak $V_{m,ph} = \sqrt{2} V_{ph,rms}$:
\begin{equation}
v_{an}(t) = \sqrt{2} V_{ph,rms}\sin(\omega t), \quad v_{bn}(t) = \sqrt{2} V_{ph,rms}\sin(\omega t - 120^\circ), \quad v_{cn}(t) = \sqrt{2} V_{ph,rms}\sin(\omega t + 120^\circ)
\end{equation}
Each diode/thyristor conducts for $120^\circ$ ($2\pi/3$ radians) when its anode has the highest phase voltage.

\paragraph{1. Uncontrolled Diode Rectifier:}
Integrating over the conduction window $[\pi/6, 5\pi/6]$:
\begin{equation}
V_{dc} = \frac{3}{2\pi}\int_{\pi/6}^{5\pi/6} V_{m,ph}\sin(\omega t)\,d(\omega t) = \frac{3\sqrt{3} V_{m,ph}}{2\pi} = \frac{3\sqrt{6}}{2\pi} V_{ph,rms} \approx 0.8270 V_{m,ph} \approx 1.1695 V_{ph,rms}
\end{equation}
\begin{align}
V_{rms} &= V_{m,ph}\sqrt{\frac{3}{2\pi}\left(\frac{\pi}{3} + \frac{\sqrt{3}}{4}\right)} \approx 0.8407 V_{m,ph} \approx 1.1889 V_{ph,rms} \\
FF &= \frac{0.8407}{0.8270} \approx 1.0166 \\
RF &= \sqrt{1.0166^2 - 1} \approx 0.183 \quad (18.3\%) \\
\text{PIV} &= \sqrt{3} V_{m,ph} = V_{m,LL} = \sqrt{2} V_{LL,rms} = \sqrt{6} V_{ph,rms} \\
f_{\mathrm{ripple}} &= 3 f_{\mathrm{line}} \quad (150\text{ Hz for 50 Hz; } 180\text{ Hz for 60 Hz})
\end{align}

\paragraph{2. Controlled (SCR) Rectifier with Continuous Current:}
With firing angle $\alpha$ measured from the natural commutation point ($\omega t = 30^\circ$):
\begin{equation}
V_{dc} = \frac{3\sqrt{3} V_{m,ph}}{2\pi}\cos\alpha \approx 1.1695 V_{ph,rms}\cos\alpha
\end{equation}

\subsubsection{Three-Phase Full-Wave Bridge Rectifier (6-Pulse)}
The standard industrial 6-pulse bridge consists of 6 diodes/thyristors:
\begin{itemize}
  \item Upper group ($D_1, D_3, D_5$): Cathodes tied to positive DC rail; the diode connected to the phase with the \textbf{most positive} voltage conducts.
  \item Lower group ($D_4, D_6, D_2$): Anodes tied to negative DC rail; the diode connected to the phase with the \textbf{most negative} voltage conducts.
  \item At any instant, exactly one diode from the upper group and one diode from the lower group conduct. The DC output voltage is the upper envelope of line-to-line voltages:
  \begin{equation}
  v_d(t) = v_{\mathrm{upper}}(t) - v_{\mathrm{lower}}(t) = v_{LL}(t)
  \end{equation}
  \item Conduction period for each device: $120^\circ$ ($2\pi/3$). Six distinct conduction pairs per line cycle ($60^\circ$ per pulse).
\end{itemize}

\paragraph{1. Diode Bridge Average DC Voltage:}
Centering the line-to-line voltage peak $V_{m,LL} \cos(\omega t)$ over $[-\pi/6, \pi/6]$:
\begin{equation}
V_{dc} = \frac{6}{2\pi}\int_{-\pi/6}^{\pi/6} V_{m,LL}\cos(\omega t)\,d(\omega t) = \frac{3 V_{m,LL}}{\pi}\left[\sin(\pi/6) - \sin(-\pi/6)\right] = \frac{3 V_{m,LL}}{\pi}
\end{equation}
In terms of RMS line-to-line and phase voltages ($V_{m,LL} = \sqrt{2} V_{LL,rms} = \sqrt{6} V_{ph,rms}$):
\begin{equation}
V_{dc} = \frac{3\sqrt{2}}{\pi} V_{LL,rms} = \frac{3\sqrt{6}}{\pi} V_{ph,rms} \approx 1.3505 V_{LL,rms} \approx 2.3391 V_{ph,rms}
\end{equation}

\paragraph{2. RMS Voltage, Ripple, and Device Stresses:}
\begin{align}
V_{rms} &= V_{LL,rms}\sqrt{1 + \frac{3\sqrt{3}}{2\pi}} \approx 1.3516 V_{LL,rms} \approx 2.3411 V_{ph,rms} \\
FF &= \frac{V_{rms}}{V_{dc}} = \frac{1.3516}{1.3505} \approx 1.0009 \\
RF &= \sqrt{FF^2 - 1} = \sqrt{1.0009^2 - 1} \approx 0.042 \quad (4.20\%) \\
\text{PIV} &= V_{m,LL} = \sqrt{2} V_{LL,rms} = \sqrt{6} V_{ph,rms} \\
f_{\mathrm{ripple}} &= 6 f_{\mathrm{line}} \quad (300\text{ Hz for 50 Hz; } 360\text{ Hz for 60 Hz})
\end{align}
\paragraph{3. Device Currents (for ripple-free DC current $I_d$):}
Each diode carries $I_d$ for $1/3$ of the fundamental cycle ($120^\circ$):
\begin{align}
I_{D,\mathrm{avg}} &= \frac{1}{2\pi}\int_0^{2\pi/3} I_d\,d(\omega t) = \frac{I_d}{3} \approx 0.3333 I_d \\
I_{D,\mathrm{rms}} &= \sqrt{\frac{1}{2\pi}\int_0^{2\pi/3} I_d^2\,d(\omega t)} = \frac{I_d}{\sqrt{3}} \approx 0.5774 I_d
\end{align}
The AC line current is a quasi-square wave with conduction duration $120^\circ$ positive, $60^\circ$ zero, $120^\circ$ negative:
\begin{equation}
I_{s,\mathrm{rms}} = \sqrt{\frac{2}{3}} I_d \approx 0.8165 I_d, \qquad I_{s1,\mathrm{rms}} = \frac{\sqrt{6}}{\pi} I_d \approx 0.7797 I_d
\end{equation}
\begin{equation}
DF = \frac{I_{s1,\mathrm{rms}}}{I_{s,\mathrm{rms}}} = \frac{\sqrt{6}/\pi}{\sqrt{2/3}} = \frac{3}{\pi} \approx 0.9549 \quad \implies \quad THD = \sqrt{\left(\frac{\pi}{3}\right)^2 - 1} \approx 0.3108 \quad (31.08\%)
\end{equation}

\subsubsection{Three-Phase Controlled SCR Bridge Rectifier (6-Pulse)}
Replacing the 6 diodes with thyristors:
\begin{equation}
V_{dc} = \frac{3 V_{m,LL}}{\pi}\cos\alpha = \frac{3\sqrt{2}}{\pi} V_{LL,rms}\cos\alpha \approx 1.3505 V_{LL,rms}\cos\alpha
\end{equation}
\begin{itemize}
  \item Rectifier operation: $0^\circ \le \alpha \le 90^\circ$ ($V_{dc} \ge 0$).
  \item Inverter operation: $90^\circ < \alpha \le 180^\circ - \gamma$ ($V_{dc} < 0$, DC power fed back to 3-phase grid).
\end{itemize}

---

\subsection{Active Boost Power Factor Correction (PFC)}

\subsubsection{Motivation and Harmonic Regulations}
Conventional diode bridge rectifiers followed by bulk reservoir capacitors draw short, peaky current pulses ($THD > 80\%$, $PF \approx 0.55 - 0.65$). This leads to:
\begin{itemize}
  \item Excessive line currents and $I^2 R$ distribution losses.
  \item Triplen harmonic currents ($3\mathrm{rd}, 9\mathrm{th}, 15\mathrm{th}, \dots$) accumulating in the neutral conductor of 3-phase systems, causing neutral overheating.
  \item Voltage flat-topping and electromagnetic interference (EMI).
  \item Violation of international emission standards: \textbf{IEC 61000-3-2} and \textbf{IEEE 519}.
\end{itemize}

\subsubsection{Boost PFC Topology and Dual-Loop Control}
An active Boost PFC preregulator consists of an uncontrolled single-phase diode bridge followed by a high-frequency Boost DC-DC converter ($L_{pfc}, S, D_{pfc}, C_o$).
\begin{itemize}
  \item The rectified input voltage is $|v_{ac}(t)| = \sqrt{2} V_{in}|\sin(\omega t)|$.
  \item \textbf{Inner Current Loop (Fast, bandwidth $\approx 10 - 20\text{ kHz}$):} Senses the instantaneous inductor current $i_L(t)$ and forces its high-frequency cycle-average $\langle i_L(t) \rangle_{T_s}$ to track the rectified AC voltage reference $|v_{ac}(t)|$. Consequently, the AC grid current $i_{ac}(t)$ is forced to be a pure sine wave in phase with $v_{ac}(t)$:
  \begin{equation}
  i_{ac}(t) = \sqrt{2} I_{s,\mathrm{rms}}\sin(\omega t) \implies \phi_1 = 0 \implies DPF = 1.0, \quad THD \approx 0 \implies DF \approx 1.0 \implies PF \approx 1.0
  \end{equation}
  \item \textbf{Outer Voltage Loop (Slow, bandwidth $\approx 10 - 20\text{ Hz}$):} Senses the DC output voltage $V_o$ and regulates it to the desired DC target. Its bandwidth must be strictly below $2 f_{\mathrm{line}}$ ($100\text{ Hz}$ or $120\text{ Hz}$) so it does not respond to or distort the inherent double-line-frequency ripple.
\end{itemize}

\subsubsection{DC Output Voltage Requirement}
Because a boost converter requires $V_o > v_{in}(t)$ at all times for proper current control:
\begin{equation}
V_o > \hat{V}_{in,\max} = \sqrt{2} V_{in,\mathrm{rms},\max}
\end{equation}
For universal AC line input ($85 - 265\,\si{\volt}_{\mathrm{rms}}$):
\begin{equation}
\hat{V}_{in,\max} = \sqrt{2} \times 265\,\si{\volt} \approx 374.8\,\si{\volt} \quad \implies \quad \text{Standard DC bus voltage is chosen as } V_o = 380\,\si{\volt} - 400\,\si{\volt}
\end{equation}

\subsubsection{Analytical Derivation of PFC Inductor Sizing ($L_{pfc}$)}
In Continuous Conduction Mode (CCM), the instantaneous duty ratio $d(t)$ of the boost switch varies over the line cycle:
\begin{equation}
d(t) = 1 - \frac{|v_{ac}(t)|}{V_o} = 1 - \frac{\sqrt{2} V_{in}|\sin(\omega t)|}{V_o}
\end{equation}
During the switch on-time $t_{\mathrm{on}}(t) = d(t) T_s = d(t)/f_s$, the inductor current ripple is:
\begin{equation}
\Delta i_L(t) = \frac{|v_{ac}(t)| d(t)}{L_{pfc} f_s} = \frac{|v_{ac}(t)|\left(V_o - |v_{ac}(t)|\right)}{L_{pfc} f_s V_o}
\end{equation}
Evaluating the ripple at the peak of the AC input voltage where $|v_{ac}| = \sqrt{2} V_{in}$:
\begin{equation}
\Delta I_L = \frac{\sqrt{2} V_{in}\left(V_o - \sqrt{2} V_{in}\right)}{L_{pfc} f_s V_o}
\end{equation}
The peak input current at minimum input voltage $V_{in}$ and rated power $P_o$ (assuming conversion efficiency $\eta$):
\begin{equation}
P_o = \eta V_{in} I_{in,\mathrm{rms}} = \eta \frac{\sqrt{2} V_{in}}{2} \hat{I}_{in} \implies \hat{I}_{in} = \frac{\sqrt{2} P_o}{\eta V_{in}}
\end{equation}
Expressing the ripple in terms of the allowable peak-to-peak ripple current $\Delta I_L$ yields the sizing equation:
\begin{equation}
\label{eq:pfc_L_sizing}
L_{pfc} \ge \frac{V_{in}^2 \left(V_o - \sqrt{2} V_{in}\right)}{f_s \Delta I_L V_o P_o}
\end{equation}
where $V_{in}$ is the nominal/minimum RMS line voltage, $V_o$ is the regulated DC output, $f_s$ is the switching frequency, and $\Delta I_L$ is the specified ripple current (typically chosen as $20\% - 40\%$ of peak input current $\hat{I}_{in}$).

\subsubsection{Analytical Derivation of Bulk Output Capacitor Sizing ($C_o$)}
The instantaneous AC input power under unity power factor is:
\begin{equation}
p_{in}(t) = v_{ac}(t) i_{ac}(t) = \left[\sqrt{2} V_{in}\sin(\omega t)\right]\left[\sqrt{2} I_{in}\sin(\omega t)\right] = 2 P_o \sin^2(\omega t) = P_o - P_o \cos(2\omega t)
\end{equation}
Meanwhile, the DC load draws constant power $P_o$. By energy conservation, the output bulk capacitor must absorb/supply the fluctuating power component at double the line frequency ($2\omega = 2\omega_{\mathrm{line}} = 4\pi f_{\mathrm{line}}$):
\begin{equation}
p_C(t) = p_{in}(t) - P_o = -P_o \cos(2\omega_{\mathrm{line}} t)
\end{equation}
The instantaneous capacitor current is:
\begin{equation}
i_C(t) = \frac{p_C(t)}{V_o} = -\frac{P_o}{V_o}\cos(2\omega_{\mathrm{line}} t)
\end{equation}
Integrating to obtain the double-line-frequency ripple voltage across $C_o$:
\begin{equation}
\tilde{v}_o(t) = \frac{1}{C_o}\int i_C(t)\,dt = -\frac{P_o}{2\omega_{\mathrm{line}} V_o C_o}\sin(2\omega_{\mathrm{line}} t)
\end{equation}
The peak ripple amplitude $\Delta V_o$ (or half of peak-to-peak ripple $\Delta V_{o,pp} = 2\Delta V_o$) is:
\begin{equation}
\Delta V_o = \frac{P_o}{2\omega_{\mathrm{line}} V_o C_o} = \frac{P_o}{4\pi f_{\mathrm{line}} V_o C_o}
\end{equation}
Rearranging gives the required bulk capacitance for a specified voltage ripple limit $\Delta V_o$:
\begin{equation}
\label{eq:pfc_Co_sizing}
C_o \ge \frac{P_o}{2\omega_{\mathrm{line}} V_o \Delta V_o} = \frac{P_o}{4\pi f_{\mathrm{line}} V_o \Delta V_o}
\end{equation}
\paragraph{Hold-Up Time Criterion:}
In telecommunication and server power supplies, $C_o$ must also sustain full load power $P_o$ during a complete AC line dropout of duration $t_{\mathrm{hold}}$ (typically $10 - 20\,\si{\milli\second}$), while voltage drops from $V_o$ to the minimum permissible operating limit $V_{o,\min}$:
\begin{equation}
\Delta E_C = \frac{1}{2} C_o \left(V_o^2 - V_{o,\min}^2\right) \ge P_o t_{\mathrm{hold}} \implies C_o \ge \frac{2 P_o t_{\mathrm{hold}}}{V_o^2 - V_{o,\min}^2}
\end{equation}
The final design selects $C_o = \max\left(C_{o,\mathrm{ripple}},\, C_{o,\mathrm{holdup}}\right)$.

---

\subsection{Exam Recipes \& Step-by-Step Problem Solving}

\subsubsection{Recipe 1: Calculating DPF, DF, THD, and True Power Factor}
\noindent\textbf{Context:} Given an AC load or rectifier current harmonic spectrum: fundamental RMS $I_{s1}$, harmonic RMS values $I_{s3}, I_{s5}, I_{s7}, \dots$, and phase displacement $\phi_1$.

\paragraph{Step-by-Step Execution Algorithm:}
\begin{enumerate}
  \item \textbf{Calculate Total RMS Current $I_{s,\mathrm{rms}}$:}
  \[
  I_{s,\mathrm{rms}} = \sqrt{I_{s1}^2 + I_{s3}^2 + I_{s5}^2 + I_{s7}^2 + \cdots}
  \]
  \item \textbf{Calculate Distortion Factor $DF$:}
  \[
  DF = \frac{I_{s1}}{I_{s,\mathrm{rms}}}
  \]
  \item \textbf{Calculate Total Harmonic Distortion $THD$:}
  \[
  THD = \frac{\sqrt{I_{s,\mathrm{rms}}^2 - I_{s1}^2}}{I_{s1}} = \sqrt{\left(\frac{1}{DF}\right)^2 - 1} = \frac{\sqrt{\sum_{n=2}^\infty I_{sn}^2}}{I_{s1}}
  \]
  \item \textbf{Calculate Displacement Power Factor $DPF$:}
  \[
  DPF = \cos\phi_1
  \]
  \item \textbf{Calculate True Power Factor $PF$:}
  \[
  PF = \frac{P}{S} = DF \times DPF = \frac{\cos\phi_1}{\sqrt{1 + THD^2}}
  \]
\end{enumerate}

\paragraph{Worked Numerical Example:}
A non-linear industrial load connected to a $230\,\si{\volt}_{\mathrm{rms}}, 50\,\si{\hertz}$ supply draws fundamental current $I_{s1} = 20.0\,\si{\ampere}$ lagging by $\phi_1 = 30^\circ$, 3rd harmonic $I_{s3} = 8.0\,\si{\ampere}$, 5th harmonic $I_{s5} = 4.0\,\si{\ampere}$, and 7th harmonic $I_{s7} = 2.0\,\si{\ampere}$.
\begin{itemize}
  \item \textbf{Step 1:} $I_{s,\mathrm{rms}} = \sqrt{20.0^2 + 8.0^2 + 4.0^2 + 2.0^2} = \sqrt{400 + 64 + 16 + 4} = \sqrt{484} = 22.0\,\si{\ampere}$.
  \item \textbf{Step 2:} $DF = \dfrac{20.0}{22.0} = \dfrac{10}{11} \approx 0.9091$.
  \item \textbf{Step 3:} $THD = \dfrac{\sqrt{22^2 - 20^2}}{20} = \dfrac{\sqrt{484 - 400}}{20} = \dfrac{\sqrt{84}}{20} = \dfrac{9.165}{20} \approx 0.4583 \quad (45.83\%)$.
  \item \textbf{Step 4:} $DPF = \cos(30^\circ) = \dfrac{\sqrt{3}}{2} \approx 0.8660$.
  \item \textbf{Step 5:} $PF = DF \times DPF = 0.9091 \times 0.8660 \approx 0.7873$.
  \item \textbf{Active Power:} $P = V_s I_{s1}\cos\phi_1 = 230 \times 20 \times 0.8660 = 3983.6\,\si{\watt}$.
  \item \textbf{Apparent Power:} $S = V_s I_{s,\mathrm{rms}} = 230 \times 22.0 = 5060.0\,\si{\volt\ampere}$.
  \item \textbf{Check:} $PF = \dfrac{P}{S} = \dfrac{3983.6}{5060.0} = 0.7873$ (Exact match!).
\end{itemize}

---

\subsubsection{Recipe 2: Determining Thyristor Firing Angle $\alpha$ and Firing Delay $t_\alpha$}
\noindent\textbf{Context:} Given an AC supply voltage and a controlled rectifier topology, determine the firing angle $\alpha$ required to achieve a specified DC output voltage $V_{dc}$ or power, and check for inversion mode.

\paragraph{Step-by-Step Execution Algorithm:}
\begin{enumerate}
  \item \textbf{Identify Circuit Configuration and Load Nature:}
  \begin{itemize}
    \item 1-$\Phi$ Full Bridge (continuous inductive current): $V_{dc} = \dfrac{2\sqrt{2} V_s}{\pi}\cos\alpha \approx 0.9003 V_s \cos\alpha$.
    \item 1-$\Phi$ Full Bridge (pure resistive load): $V_{dc} = \dfrac{\sqrt{2} V_s}{\pi}(1 + \cos\alpha) \approx 0.4502 V_s (1 + \cos\alpha)$.
    \item 3-$\Phi$ Full Bridge (6-pulse, continuous current): $V_{dc} = \dfrac{3\sqrt{2} V_{LL,rms}}{\pi}\cos\alpha \approx 1.3505 V_{LL,rms} \cos\alpha$.
  \end{itemize}
  \item \textbf{Solve for $\cos\alpha$:}
  \[
  \cos\alpha = \frac{V_{dc}}{V_{dc0}}
  \]
  where $V_{dc0}$ is the maximum uncontrolled DC voltage at $\alpha = 0^\circ$. Ensure $|\cos\alpha| \le 1$.
  \item \textbf{Compute Firing Angle $\alpha$:}
  \[
  \alpha = \arccos\left(\frac{V_{dc}}{V_{dc0}}\right)
  \]
  \begin{itemize}
    \item If $V_{dc} > 0 \implies 0^\circ \le \alpha < 90^\circ$ (\textbf{Rectifier Mode}).
    \item If $V_{dc} < 0 \implies 90^\circ < \alpha < 180^\circ$ (\textbf{Inverter Mode}; verify back-EMF $E > |V_{dc}|$ and $\alpha < 180^\circ - \gamma$).
  \end{itemize}
  \item \textbf{Compute Firing Delay Time $t_\alpha$:}
  \[
  t_\alpha = \frac{\alpha \text{ [rad]}}{2\pi f_{\mathrm{line}}} = \frac{\alpha \text{ [deg]}}{360^\circ \times f_{\mathrm{line}}}
  \]
\end{enumerate}

\paragraph{Worked Numerical Example:}
A 3-phase 6-pulse controlled SCR bridge is fed from a $400\,\si{\volt}_{\mathrm{rms}}$ line-to-line, $50\,\si{\hertz}$ AC grid supplying a large inductive DC motor load.
\begin{enumerate}
  \item[(a)] \textbf{Find $\alpha$ to deliver $V_{dc} = +300\,\si{\volt}$ (Motoring / Rectifier mode):}
  \begin{align*}
  V_{dc0} &= \frac{3\sqrt{2}}{\pi} \times 400 = 1.3505 \times 400 = 540.19\,\si{\volt} \\
  \cos\alpha &= \frac{300}{540.19} = 0.5554 \implies \alpha = \arccos(0.5554) = 56.26^\circ \approx 0.982\,\si{\radian} \\
  t_\alpha &= \frac{56.26^\circ}{360^\circ \times 50\,\si{\hertz}} = \frac{56.26}{18000} = 3.126\,\si{\milli\second}
  \end{align*}
  \item[(b)] \textbf{Find $\alpha$ for regenerative braking (Inverter mode) with $V_{dc} = -300\,\si{\volt}$:}
  \begin{align*}
  \cos\alpha &= \frac{-300}{540.19} = -0.5554 \implies \alpha = 180^\circ - 56.26^\circ = 123.74^\circ \approx 2.160\,\si{\radian} \\
  t_\alpha &= \frac{123.74^\circ}{18000} = 6.874\,\si{\milli\second}
  \end{align*}
  \textit{Commutation check:} $180^\circ - 123.74^\circ = 56.26^\circ > \gamma_{\min} \approx 20^\circ$, ensuring safe margin against commutation failure.
\end{enumerate}

---

\subsubsection{Recipe 3: Sizing PFC Boost Inductor $L_{pfc}$ and Bulk Capacitor $C_o$}
\noindent\textbf{Context:} Design an active Boost PFC converter operating from AC grid supply to produce a regulated DC bus voltage.

\paragraph{Given Design Specifications:}
\begin{itemize}
  \item Nominal/minimum input voltage: $V_{in,\mathrm{rms}}$ at line frequency $f_{\mathrm{line}}$.
  \item Regulated DC output voltage: $V_o$ (must satisfy $V_o > \sqrt{2} V_{in,\max}$).
  \item Rated DC output power: $P_o$ (and efficiency $\eta$).
  \item Switching frequency: $f_s$.
  \item Allowable peak-to-peak inductor current ripple fraction: $r = \Delta I_L / \hat{I}_{in}$ (typically $20\% - 40\%$).
  \item Maximum allowable peak-to-peak output voltage ripple: $\Delta V_{o,pp} = 2\Delta V_o$.
  \item Required hold-up time: $t_{\mathrm{hold}}$ with minimum output voltage $V_{o,\min}$.
\end{itemize}

\paragraph{Step-by-Step Execution Algorithm:}
\begin{enumerate}
  \item \textbf{Verify Boost Voltage Condition:} Confirm $V_o > \sqrt{2} V_{in,\max}$.
  \item \textbf{Determine Peak Input Current $\hat{I}_{in}$ and Inductor Ripple $\Delta I_L$:}
  \[
  \hat{I}_{in} = \frac{\sqrt{2} P_o}{\eta V_{in,\mathrm{rms}}}, \qquad \Delta I_L = r \hat{I}_{in}
  \]
  \item \textbf{Calculate Minimum Boost Inductance $L_{pfc}$:}
  \[
  L_{pfc} \ge \frac{V_{in}^2 (V_o - \sqrt{2} V_{in})}{f_s \Delta I_L V_o P_o}
  \]
  \item \textbf{Calculate Minimum Bulk Capacitance for $2 f_{\mathrm{line}}$ Ripple:}
  \[
  C_{o,\mathrm{ripple}} \ge \frac{P_o}{2\omega_{\mathrm{line}} V_o \Delta V_o} = \frac{P_o}{4\pi f_{\mathrm{line}} V_o \Delta V_o} = \frac{P_o}{2\pi f_{\mathrm{line}} V_o \Delta V_{o,pp}}
  \]
  \item \textbf{Calculate Minimum Capacitance for Hold-Up Time:}
  \[
  C_{o,\mathrm{holdup}} \ge \frac{2 P_o t_{\mathrm{hold}}}{V_o^2 - V_{o,\min}^2}
  \]
  \item \textbf{Select Final Capacitance:} $C_o = \max\left(C_{o,\mathrm{ripple}},\, C_{o,\mathrm{holdup}}\right)$.
\end{enumerate}

\paragraph{Worked Numerical Example:}
Design a Boost PFC preregulator with specifications:
\begin{itemize}
  \item AC input: $V_{in} = 230\,\si{\volt}_{\mathrm{rms}}, f_{\mathrm{line}} = 50\,\si{\hertz}$.
  \item Output voltage: $V_o = 400\,\si{\volt}$.
  \item Rated output power: $P_o = 1000\,\si{\watt}$ ($1.0\,\si{\kilo\watt}$), assumed efficiency $\eta = 0.95$.
  \item Switching frequency: $f_s = 100\,\si{\kilo\hertz}$ ($10^5\,\si{\hertz}$).
  \item Inductor current ripple ratio: $r = 25\%$ of peak input current.
  \item Output voltage ripple limit: $\Delta V_o \le 5\,\si{\volt}$ peak ($\Delta V_{o,pp} \le 10\,\si{\volt}$).
  \item Hold-up time: $t_{\mathrm{hold}} = 10\,\si{\milli\second}$ with $V_{o,\min} = 340\,\si{\volt}$.
\end{itemize}

\paragraph{Solution:}
\begin{itemize}
  \item \textbf{Step 1 (Boost Condition):} $\hat{V}_{in} = \sqrt{2} \times 230 = 325.27\,\si{\volt} < 400\,\si{\volt}$ (Satisfied).
  \item \textbf{Step 2 (Currents):}
  \[
  \hat{I}_{in} = \frac{\sqrt{2} \times 1000}{0.95 \times 230} = \frac{1414.2}{218.5} = 6.472\,\si{\ampere}
  \]
  \[
  \Delta I_L = 0.25 \times 6.472\,\si{\ampere} = 1.618\,\si{\ampere}
  \]
  \item \textbf{Step 3 (Inductor Sizing via Eq.~\ref{eq:pfc_L_sizing}):}
  \begin{align*}
  L_{pfc} &\ge \frac{(230)^2 \times (400 - \sqrt{2} \times 230)}{100\times 10^3 \times 1.618 \times 400 \times 1000} = \frac{52900 \times (400 - 325.27)}{6.472 \times 10^{10}} \\
  &= \frac{52900 \times 74.73}{6.472 \times 10^{10}} = \frac{3.953 \times 10^6}{6.472 \times 10^{10}} = 6.108 \times 10^{-5}\,\si{\henry}
  \end{align*}
  \[
  L_{pfc} \ge 61.1\,\si{\micro\henry} \quad \implies \quad \text{Select standard } L_{pfc} \ge 65\,\si{\micro\henry} \text{ or } 100\,\si{\micro\henry}.
  \]
  \item \textbf{Step 4 (Capacitor Sizing for $100\,\si{\hertz}$ Ripple via Eq.~\ref{eq:pfc_Co_sizing}):}
  \[
  \omega_{\mathrm{line}} = 2\pi (50) = 100\pi \approx 314.16\,\si{\radian\per\second}
  \]
  \[
  C_{o,\mathrm{ripple}} \ge \frac{1000}{2 \times (100\pi) \times 400 \times 5} = \frac{1000}{400000\pi} = \frac{1}{400\pi} \approx 7.958 \times 10^{-4}\,\si{\farad} = 795.8\,\si{\micro\farad}
  \]
  \item \textbf{Step 5 (Hold-Up Capacitance):}
  \[
  C_{o,\mathrm{holdup}} \ge \frac{2 \times 1000 \times 10 \times 10^{-3}}{(400)^2 - (340)^2} = \frac{20}{160000 - 115600} = \frac{20}{44400} = 4.505 \times 10^{-4}\,\si{\farad} = 450.5\,\si{\micro\farad}
  \]
  \item \textbf{Step 6 (Final Selection):}
  \[
  C_o \ge \max(795.8\,\si{\micro\farad},\, 450.5\,\si{\micro\farad}) = 795.8\,\si{\micro\farad} \quad \implies \quad \text{Select } C_o \ge 820\,\si{\micro\farad} \text{ or } 1000\,\si{\micro\farad} \, (450\,\si{\volt}).
  \]
\end{itemize}

---

\subsection{Common Exam Pitfalls \& Traps}
\begin{center}
\fbox{\parbox{0.92\textwidth}{
\begin{description}
  \item[Trap 1: Line-to-Line vs Phase Voltage in 3-Phase Equations.] In three-phase bridges, the DC voltage is $V_{dc} = \dfrac{3 V_{m,LL}}{\pi} = \dfrac{3\sqrt{2}}{\pi} V_{LL,rms} \approx 1.35 V_{LL,rms}$. Do NOT plug phase voltage into this equation without multiplying by $\sqrt{3}$! ($V_{LL,rms} = \sqrt{3} V_{ph,rms}$).
  \item[Trap 2: Conflating Peak $V_m$ with RMS $V_s$.] $V_m = \sqrt{2} V_s$. In half-wave rectifiers, $V_{dc} = V_m/\pi \approx 0.318 V_m$, which equals $0.450 V_s$. Check whether the exam prompt provides peak or RMS voltage before multiplying!
  \item[Trap 3: Assuming $DPF = 1$ implies $PF = 1$.] Power factor correction must address both displacement ($\cos\phi_1$) and harmonic distortion ($THD$). If current is in phase with voltage ($\phi_1 = 0$) but contains $50\%$ THD, $PF = 1/\sqrt{1 + 0.5^2} = 0.894 \ne 1.0$.
  \item[Trap 4: Inverter Mode Conditions ($\alpha > 90^\circ$).] A thyristor bridge operates as an inverter ONLY if: (1) current is continuous, (2) an external active DC source exists ($E > |V_{dc}|$), and (3) polarity is aligned. For a simple passive resistive load, current ceases when voltage reverses, so $V_{dc}$ cannot be negative and the bridge cannot invert!
  \item[Trap 5: PFC Bulk Capacitor Ripple Frequency.] The ripple voltage across $C_o$ occurs at \textbf{twice the line frequency} ($2 f_{\mathrm{line}} = 100\,\si{\hertz}$ or $120\,\si{\hertz}$), NOT the switching frequency ($f_s = 100\,\si{\kilo\hertz}$). Sizing $C_o$ using $f_s$ produces an error of three orders of magnitude!
  \item[Trap 6: Boost PFC Voltage Violation ($V_o < \sqrt{2} V_{in}$).] The Boost PFC converter cannot step down. If $V_o \le \sqrt{2} V_{in}$, the boost diode remains forward-biased, uncontrollable current surges occur, and PFC control fails completely.
\end{description}
}}
\end{center}
